\chapter{Marchés d'options en Europe et en Asie}\label{ch:m1:options-markets-europe-asia}

En 2024, plus de quatre contrats à terme et d'options sur cinq échangés sur
les bourses du monde l'étaient en Asie-Pacifique~; un an plus tard, moins
de deux sur trois, et le total mondial avait reculé de quatre dixièmes alors que
le volume des contrats à terme augmentait. Aucune crise n'avait eu lieu. Un régulateur avait jugé qu'un contrat
que des particuliers achetaient en nombre énorme était trop petit et expirait
trop souvent à échéance, et il en avait réécrit la spécification. Douze ans plus tôt, un autre
régulateur avait fait de même avec ce qui était alors l'option la plus échangée
au monde, pour la même raison. Les marchés d'options hors des
États-Unis sont nationaux, chacun avec une bourse dominante, sa propre
conception de contrat, sa propre population d'investisseurs particuliers et sa
propre vision de ces deux éléments par son régulateur. Ce chapitre donne les outils pour les comparer, et le premier
consiste à refuser de compter les contrats.

\section{Trois façons de mesurer un marché}

\begin{definition}[Taille de lot, notionnel et volume en primes]\label{def:m1:options-markets-europe-asia:turnover}
La \emph{taille de lot}\index{taille de lot (dérivés)} d'un produit dérivé est le
nombre d'unités de sous-jacent par contrat~: c'est un autre nom du multiplicateur,
employé lorsque les contrats sont spécifiés en actions ou en unités d'indice par lot.
Le \emph{volume en notionnel}\index{volume en notionnel} est égal au
nombre de contrats multiplié par la taille de lot et par le prix du sous-jacent. Le \emph{volume en primes}\index{volume en primes}
est, pour les options, égal au nombre de contrats multiplié par la taille de lot et
par le prix de l'option~: c'est la somme d'argent qui a effectivement changé de mains.
\end{definition}

Les trois mesures répondent à des questions différentes. Le nombre de contrats mesure les messages
et les frais de bourse. Le notionnel mesure l'exposition qui a changé de mains, et
il la surestime pour les options, dont la plupart sont très hors de la monnaie. Les primes
mesurent la somme d'argent en jeu et le fonds de revenus disponible pour les \omterm{def:m1:what-a-trading-firm-does:market-maker}{teneurs de marché},
dont la fourchette ne représente qu'une fraction de la prime.
La \Cref{fig:m1:options-markets-europe-asia:rankings} classe quatre marchés fictifs,
aux proportions réalistes, selon ces trois mesures.

\begin{omfigure}
\begin{tikzpicture}
\begin{axis}[width=11cm,height=5.2cm,ybar,bar width=10pt,
  ylabel={part (\%)},
  symbolic x coords={M1,M2,M3,M4},xtick=data,
  xticklabels={{M1\\hebdo. petit lot},{M2\\grand indice},{M3\\indice moyen},{M4\\futures indiciels}},
  tick label style={font=\small},label style={font=\small},xticklabel style={font=\footnotesize,align=center},
  ymin=0,ymax=112,ymajorgrids,grid style={gray!25},enlarge x limits=0.15,
  nodes near coords,nodes near coords style={font=\scriptsize,/pgf/number format/.cd,fixed,precision=0},
  legend columns=3,legend style={font=\footnotesize,at={(0.5,-0.3)},anchor=north,draw=none,fill=none,/tikz/every even column/.append style={column sep=8pt}},
  table/col sep=comma]
\addplot[fill=omExo!50,draw=omExo] table[x=market,y=contracts]{figdata/markets-1/27-options-markets-europe-asia/rankings.csv};
\addplot[fill=omDef!70,draw=omDef] table[x=market,y=notional]{figdata/markets-1/27-options-markets-europe-asia/rankings.csv};
\addplot[fill=omThm!60,draw=omThm] table[x=market,y=premium]{figdata/markets-1/27-options-markets-europe-asia/rankings.csv};
\legend{par contrats,par notionnel,par primes}
\end{axis}
\end{tikzpicture}

\omcaption{Quatre marchés fictifs. M1, des \omterm{def:m1:zero-day-options-and-retail-flow:zero}{options hebdomadaires} sur indice à petit lot,
représente 98~\% des contrats, 48~\% du notionnel et 38~\% des primes~; M2, une option sur
indice à gros contrat, représente moins de 1~\% des contrats et 55~\% des
primes. Données~: le tutoriel.}\label{fig:m1:options-markets-europe-asia:rankings}
\end{omfigure}

\begin{omfigure}
\begin{tikzpicture}
\begin{axis}[width=10.5cm,height=4.8cm,xbar,bar width=6pt,
  xlabel={contrats (milliards)},
  symbolic y coords={Other,Europe,Latin America,North America,Asia-Pacific},ytick={Other,Europe,Latin America,North America,Asia-Pacific},yticklabels={{Autres},{Europe},{Amérique latine},{Amérique du Nord},{Asie-Pacifique}},
  tick label style={font=\small},label style={font=\small},
  xmin=0,xmax=200,xmajorgrids,grid style={gray!25},enlarge y limits=0.14,
  nodes near coords,nodes near coords style={font=\scriptsize,fill=white,inner sep=1pt,/pgf/number format/.cd,fixed,fixed zerofill,precision=1},
  legend cell align=left,legend style={font=\footnotesize,at={(0.97,0.5)},anchor=east,draw=none,fill=white},reverse legend,
  table/col sep=comma]
\addplot[fill=omDef!60,draw=omDef] table[y=region,x=y2025]{figdata/markets-1/27-options-markets-europe-asia/regions.csv};
\addplot[fill=omExo!40,draw=omExo] table[y=region,x=y2024]{figdata/markets-1/27-options-markets-europe-asia/regions.csv};
\legend{2025,2024 (dérivé)}
\end{axis}
\end{tikzpicture}

\omcaption{Volume des produits dérivés cotés par région. Les chiffres de 2025 sont
ceux de l'association professionnelle~; ceux de 2024 sont dérivés de ses
variations annuelles publiées. L'Asie-Pacifique est passée de 170 à 76 milliards de
contrats, tandis que toutes les autres régions progressaient. Données~: \cref{dat:m1:index-and-single-stock-futures-worldwide:fia}.}\label{fig:m1:options-markets-europe-asia:regions}
\end{omfigure}

\section{Europe}

Les options sur actions européennes se concentrent sur deux groupes de bourses, chacun
abritant ses indices et actions nationaux, chacun compensant dans sa propre chambre
(\cref{ch:m1:option-contracts-and-exchanges}). Le produit phare est l'option sur
l'indice vedette de la zone euro~: de style européen, réglée en espèces, à \euro10 le point
comme son future, couverte avec ce future et cotée par rapport à lui. Trois
caractéristiques distinguent ces marchés du marché américain. Une part importante du
volume institutionnel se négocie de gré à gré et s'enregistre auprès de la
bourse sous forme de bloc, de sorte que l'écran n'en montre qu'une fraction de l'activité.
La participation des particuliers aux options cotées est faible~; ceux qui veulent
un effet de levier achètent des warrants et des certificats émis par les banques, un marché parallèle avec
ses propres plateformes et ses propres \omterm{def:m1:what-a-trading-firm-does:market-maker}{teneurs de marché}, les émetteurs eux-mêmes. Enfin, il
n'existe qu'un seul carnet d'ordres par produit, de sorte que le routage et les mécanismes
d'enchères du \cref{ch:m1:options-market-structure} n'ont aucun équivalent~: la structure
est celle d'une bourse de futures.

\section{Corée}

Dans les années 2000, le produit dérivé le plus échangé au monde, en nombre de contrats, était le
option sur l'indice vedette coréen~: un petit contrat, des strikes bon marché
hors de la monnaie, et une très large population de traders particuliers.
traders particuliers.

\begin{dated}{2012-06}{Le multiplicateur, multiplié par cinq}\label{dat:m1:options-markets-europe-asia:kospi}
Le 14 juin 2012, la taille du contrat de l'option sur indice a été portée de 100\,000
à 500\,000 won par point, sur la bourse d'origine et sur la bourse européenne
qui cotait un produit lié.
\end{dated}

Le changement a multiplié par cinq le prix du billet de loterie le moins cher.
Si l'activité reste inchangée, le nombre de contrats doit être divisé par cinq~; de combien
l'activité elle-même baisse est la question de l'exercice~3. C'est le
modèle de ce qui a suivi ailleurs, et un avertissement permanent pour toute société
dont l'activité dépend du flux des particuliers sur un seul marché~: la
spécification du contrat est un instrument de politique publique.

\section{Inde}

\begin{definition}[Échéance hebdomadaire et limite de position]\label{def:m1:options-markets-europe-asia:weekly}
Une \emph{échéance hebdomadaire}\index{échéance hebdomadaire} est une échéance d'option cotée chaque
semaine à un jour fixe, en plus du cycle mensuel. Une \emph{limite de
position}\index{limite de position} est un plafond fixé par une bourse ou un régulateur sur la
taille de la position qu'un intervenant ou un groupe peut détenir dans un contrat,
mesurée sur une base déclarée (brute ou nette, en contrats, en notionnel ou
ajustée du \omterm{def:m1:zero-day-options-and-retail-flow:gamma}{delta}) et contrôlée à des moments déterminés.
\end{definition}

En 2024, le marché indien des options sur indices était le plus important au monde en
nombre de contrats, d'un ordre de grandeur~: des lots de petite taille, plusieurs indices dont chacun a une
\omterm{def:m1:options-markets-europe-asia:weekly}{échéance hebdomadaire} à un jour de la semaine différent, de sorte qu'un indice arrivait à échéance
chaque jour, et des millions de particuliers négociant des options le jour de leur échéance.

\begin{dated}{2026-09}{Ce que le régulateur a constaté, et ce qu'il a fait}\label{dat:m1:options-markets-europe-asia:sebi}
\emph{L'enquête.} L'étude du régulateur des marchés financiers, publiée le 23 septembre
2024, a établi que 93~\% des particuliers intervenant sur les dérivés sur actions ont perdu de l'argent
au cours des trois exercices financiers clos en mars 2024, avec des pertes cumulées supérieures
à 1.8 billion de roupies.
\emph{Les mesures.} Sa circulaire du 1\up{er}~octobre 2024 a fixé la valeur des contrats de dérivés
sur indices à 1.5 à 2 millions de roupies lors de leur révision (à partir du 20 novembre
2024), a autorisé chaque bourse à ne proposer des \omterm{def:m1:options-markets-europe-asia:weekly}{échéances hebdomadaires} que sur un seul indice de référence,
a exigé que les primes d'options soient encaissées d'avance auprès des acheteurs, a supprimé le
avantage de marge des \omterm{def:m1:matching-algorithms-and-implied-spreads:calendar}{calendar spreads} le jour de l'échéance, a ajouté une marge de perte extrême de 2~\%
sur les options vendues le jour de l'échéance, et a introduit un suivi intrajournalier des
\omterm{def:m1:options-markets-europe-asia:weekly}{limites de position}.
\emph{La décision.} Le 3 juillet 2025, le régulateur a pris une injonction provisoire contre
un groupe étranger de trading pour compte propre, l'accusant de manipulation d'un indice bancaire les
jours d'échéance entre janvier 2023 et 2025, et a saisi 48.4 milliards de roupies
de gains présumés. Le groupe a consigné ce montant le 14 juillet 2025 tout en
se réservant le droit de contester l'injonction. L'affaire n'était pas tranchée au
moment de la rédaction.
\end{dated}

\section{Le cas du jour d'échéance}

Mettons de côté la question de ce qu'une société a fait, à laquelle un tribunal répondra,
et considérons la structure que décrit l'injonction, car elle est générale.
Un indice compte quelques dizaines de composantes dont les marchés au comptant et à terme ont une
certaine profondeur. Les options sur l'indice, qui arrivent à échéance aujourd'hui, ont un \omterm{def:m1:options-markets-europe-asia:turnover}{volume en notionnel}
plusieurs fois supérieur à cette profondeur. Les options se règlent sur la valeur de clôture de l'indice,
calculée à partir des cours des composantes sur les dernières minutes. Dès lors, un
intervenant suffisamment grand pour faire bouger les composantes d'une fraction de pour cent
modifie le gain d'une position d'options plusieurs fois supérieure aux actions
échangées.

\begin{proposition}[Effet de levier du règlement]\label{prop:m1:options-markets-europe-asia:leverage}
Soit un intervenant détenant des options sur indices dont l'exposition agrégée équivalente à \omterm{def:m1:delta-one-instruments:deltaone}{delta-one}
$N$ à l'échéance (le notionnel des options qui finissent dans la
monnaie), et soit le déplacement de la valeur de règlement de l'indice d'une fraction
$\varepsilon$ coûtant $c(\varepsilon)$ en impact sur les actions sous-jacentes. Le
gain varie de $N\varepsilon$. Avec la règle de la racine carrée,
$c(\varepsilon) \propto \varepsilon^{3}$ pour les actions échangées afin de produire
$\varepsilon$, de sorte que, pour de petits $\varepsilon$, le gain est linéaire et le coût
cubique~: dès que $N$ est grand par rapport à la profondeur du sous-jacent, il existe
un $\varepsilon > 0$ rentable.
\end{proposition}
\begin{proof}
D'après la \cref{met:m1:the-sell-side:sqrt}, l'impact $\varepsilon \propto
\sigma\sqrt{Q/V}$, donc la quantité nécessaire est $Q \propto V\varepsilon^2/\sigma^2$
et le coût de sa négociation, quantité multipliée par l'impact moyen, est proportionnel à
$Q\varepsilon \propto \varepsilon^3$.
\end{proof}

C'est pourquoi les procédures de règlement utilisent des moyennes sur une fenêtre, pourquoi les régulateurs
surveillent les positions du jour d'échéance en cours de séance, et pourquoi la fonction conformité d'une société
doit examiner \emph{conjointement} ses livres au comptant et ses livres de dérivés~: des transactions
défendables prises isolément comme couverture ou arbitrage peuvent, combinées et avec une intention,
constituer une manipulation, et le schéma reste visible a posteriori dans exactement les
données que la société elle-même conserve. La distinction entre le fait de faire bouger un cours
comme effet secondaire d'une négociation légitime et le fait de négocier pour le faire bouger est tracée
par le droit, non par l'arithmétique, et varie selon les juridictions~; le livre~16
y revient.

\begin{omfigure}
\begin{tikzpicture}[font=\small,
  bk/.style={draw=omDef,rounded corners=2pt,minimum width=3.4cm,minimum height=1.1cm,align=center,fill=omDef!4}]
\node[bk] (und) at (0,0) {Actions composantes\\et leurs futures};
\node[bk,draw=omThm,fill=omThm!5] (opt) at (8.2,0) {Options sur indice\\arrivant à échéance aujourd'hui};
\node[bk,draw=omExo,fill=omExo!6,minimum width=3.0cm] (idx) at (4.1,2.1) {Valeur de règlement de l'indice};
\draw[-{Stealth[length=2mm]},thick,omDef] (und.north) |- node[pos=0.25,left,font=\footnotesize,align=right]{cours sur la\\fenêtre finale} (idx.west);
\draw[-{Stealth[length=2mm]},thick,omThm] (idx.east) -| node[pos=0.75,right,font=\footnotesize,align=left]{gain de chaque\\option dans la monnaie} (opt.north);
\node[font=\footnotesize,below=4pt of und] {profondeur~: $V$};
\node[font=\footnotesize,below=4pt of opt] {exposition~: $N \gg V$};
\end{tikzpicture}

\omcaption{La structure commune à tous les cas de manipulation du règlement~: un petit
marché dont les cours déterminent le gain d'un marché beaucoup plus
grand.}\label{fig:m1:options-markets-europe-asia:structure}
\end{omfigure}

\section{Autres marchés asiatiques}

Ailleurs en Asie, la situation est plus proche de celle de l'Europe~: une bourse de dérivés
par pays, des options sur indices comme produit principal, \omterm{def:m1:option-contracts-and-exchanges:style}{exercice européen} et règlement en espèces,
structure de marché de type contrats à terme. Les options sur indices du Japon et de Hong Kong
sont dominées par les institutionnels et par la couverture des produits structurés
vendus aux particuliers, dont les émetteurs sont structurellement longs ou courts
de volatilité, d'une manière qui façonne la surface locale~; les options sur indices de Taïwan
ont une large part de particuliers. Pour une société de trading, les questions du
\cref{ch:m1:asian-and-emerging-equity-markets} se posent d'abord dans chacun de ces cas~: qui peut
négocier, par l'intermédiaire de qui, avec quelles \omterm{def:m1:options-markets-europe-asia:weekly}{limites de position}, et selon quelle fiscalité.

\section{Tutoriel~: trois classements}\label{tut:m1:options-markets-europe-asia:rankings}

\begin{tutorial}
\textbf{Objectif.} Convertir les nombres de contrats en \omterm{def:m1:options-markets-europe-asia:turnover}{volumes en notionnel} et en \omterm{def:m1:options-markets-europe-asia:turnover}{volumes en primes}
dans une même devise, classer les marchés de trois façons et calculer ce que doit produire une modification
de la \omterm{def:m1:options-markets-europe-asia:turnover}{taille de lot} sur un nombre de contrats. \textbf{État final~:} les deux figures de données
de ce chapitre.
\begin{tutsteps}
\item \textbf{Un enregistrement par marché.}
\omcode{code/firm/turnover/firm_turnover.py}{9}{30}{Nombre de contrats, taille de lot, sous-jacent, prime, devise~: trois mesures du volume.}\label{lst:m1:options-markets-europe-asia:stats}
\item \textbf{Modifications de la \omterm{def:m1:options-markets-europe-asia:turnover}{taille de lot}.} Si l'activité en notionnel reste inchangée, tripler la
\omterm{def:m1:options-markets-europe-asia:turnover}{taille de lot} divise le nombre de contrats par trois~; le nombre observé ensuite indique
quelle activité a été perdue.
\omcode{code/firm/turnover/firm_turnover.py}{46}{54}{L'arithmétique d'une modification de la taille de lot, dans les deux sens.}\label{lst:m1:options-markets-europe-asia:lot}
\end{tutsteps}
\textbf{Pour aller plus loin.} Remplacer la prime moyenne par une distribution
sur les prix d'exercice et estimer le revenu potentiel d'un \omterm{def:m1:what-a-trading-firm-does:market-maker}{teneur de marché} comme la moitié de la fourchette cotée
multipliée par le \omterm{def:m1:options-markets-europe-asia:turnover}{volume en primes}, par marché.
\end{tutorial}

\section{Construction~: le normaliseur de volumes}\label{bld:m1:options-markets-europe-asia:turnover}

\begin{build}
\bfield{Rôle} Lorsque le comité de stratégie de la petite société demande
quel marché aborder ensuite, la première diapositive ne doit pas être un classement
des marchés par nombre de contrats. \bfield{Interface} \texttt{MarketStats(name, kind, contracts, multiplier,
underlying, avg\_premium, fx\_to\_usd)} avec \texttt{notional\_usd},
\texttt{premium\_usd}, \texttt{premium\_to\_notional\_bp};
\texttt{shares(markets, measure)}; \texttt{ranking(markets, measure)};
\texttt{contracts\_after\_lot\_change(\ldots)}; \texttt{activity\_kept(\ldots)}.
\bfield{Règles} Une seule devise pour toutes les comparaisons, avec le taux de change
conservé à côté des données. Les contrats à terme n'ont pas de \omterm{def:m1:options-markets-europe-asia:turnover}{volume en primes} et sont exclus,
avec une erreur, d'un classement par primes portant sur les seuls contrats à terme. Les égalités sont départagées
par nom.
\bfield{Tests d'acceptation} \texttt{code/firm/turnover/tests/}~: unités~; trois
mesures donnant trois classements~; l'arithmétique de la modification de la \omterm{def:m1:options-markets-europe-asia:turnover}{taille de lot} dans les deux sens.
\bfield{Pour aller plus loin} Notionnel ajusté du \omterm{def:m1:zero-day-options-and-retail-flow:gamma}{delta} pour les options, à partir d'une distribution
du volume par strike.
\end{build}

\begin{omsources}
\item Securities and Exchange Board of India, communiqué de presse du 23 septembre
2024 sur l'étude actualisée des traders particuliers sur les produits dérivés sur actions~;
circulaire SEBI/HO/MRD/TPD-1/P/CIR/2024/132 du 1\up{er}~octobre 2024.
\item Securities and Exchange Board of India, ordonnance provisoire du 3 juillet 2025 dans
l'affaire de manipulation d'indice~; commentaire~: ``Jane Street and the Expiry
Day Trap'', \emph{Oxford Business Law Blog}, juillet 2025.
\item Eurex, ``Eurex and Korea Exchange further expand their link'' (sur le
changement de taille de contrat de juin 2012).
\item FIA, \emph{ETD Volume --- December 2025}.
\end{omsources}

\section{Exercices}

\begin{exercise}[$\star$]\label{exo:m1:options-markets-europe-asia:1}
Une option sur indice a une \omterm{def:m1:options-markets-europe-asia:turnover}{taille de lot} de 25, l'indice vaut 22\,000 roupies et
l'option coûte 35 roupies par unité. La roupie valant 0.012~\$, calculer le notionnel
et la prime d'un contrat en dollars.
\end{exercise}

\begin{exercise}[$\star$]\label{exo:m1:options-markets-europe-asia:2}
Même question pour une option sur indice avec un multiplicateur de 100, l'indice à 5\,600~dollars
et une prime de 16. Combien de contrats de l'exercice précédent représentent la
prime d'un de ceux-ci~?
\end{exercise}

\begin{exercise}[$\star$]\label{exo:m1:options-markets-europe-asia:3}
Un contrat s'échange 90 milliards de fois par an avec une \omterm{def:m1:options-markets-europe-asia:turnover}{taille de lot} de 25. La \omterm{def:m1:options-markets-europe-asia:turnover}{taille de lot} est portée
à 75. Quel nombre de contrats faut-il attendre si l'activité en notionnel reste inchangée~? Le
nombre de contrats s'avère être de 24 milliards. Qu'est-il arrivé à l'activité~?
\end{exercise}

\begin{exercise}[$\star\star$]\label{exo:m1:options-markets-europe-asia:4}
À partir de la \cref{fig:m1:options-markets-europe-asia:regions}, calculer la part de l'Asie-Pacifique
dans le nombre mondial de contrats en 2024 et en 2025, ainsi que le total mondial de 2024.
Vérifier le total par rapport à la baisse publiée de 42.2~\%.
\end{exercise}

\begin{exercise}[$\star\star$]\label{exo:m1:options-markets-europe-asia:5}
Un \omterm{def:m1:what-a-trading-firm-does:market-maker}{teneur de marché} gagne en moyenne un quart de la fourchette cotée, et les fourchettes représentent
4~\% de la prime sur le marché M1 et 1.5~\% sur le marché M2. En utilisant des \omterm{def:m1:options-markets-europe-asia:turnover}{volumes en primes} de
945~milliards de dollars et 1\,360~milliards de dollars, estimer le revenu potentiel de chacun. Quel
marché aborder, et que faudrait-il savoir de plus~?
\end{exercise}

\begin{exercise}[$\star\star$]\label{exo:m1:options-markets-europe-asia:6}
Donner deux raisons pour lesquelles une valeur de règlement calculée comme une moyenne sur les
trente dernières minutes est plus difficile à faire bouger qu'un cours de la dernière transaction, ainsi qu'un coût
de son utilisation.
\end{exercise}

\begin{exercise}[$\star\star\star$]\label{exo:m1:options-markets-europe-asia:7}
\emph{Codage.} Avec les quatre marchés du tutoriel, afficher les trois tableaux
de parts. Puis porter la \omterm{def:m1:options-markets-europe-asia:turnover}{taille de lot} de M1 à 75 avec \texttt{activity\_kept} égal à 0.8 et
les afficher à nouveau. Quelles parts changent~?
\end{exercise}

\begin{exercise}[$\star\star\star$]\label{exo:m1:options-markets-europe-asia:8}
\emph{Trouver la faille.} Un responsable de desk~: «~Le matin, nous avons acheté des actions bancaires
parce qu'elles étaient bon marché par rapport aux contrats à terme. L'après-midi, nous les avons vendues
parce que notre limite de risque l'exigeait. Par ailleurs, notre desk d'options était vendeur de l'indice.
Chaque transaction a sa propre justification, il n'y a donc rien à expliquer.~» Que
manque-t-il à cette défense~?
\end{exercise}

\section{Problème~: compter les contrats}

\begin{problem}[{Problème du week-end --- où est l'argent~?}]\label{pb:m1:options-markets-europe-asia:1}
Une société de trading compare deux marchés d'options sur indice sur une année. M1~: 90 milliards
contrats, \omterm{def:m1:options-markets-europe-asia:turnover}{taille de lot} de 25, indice à 22\,000 roupies, prime moyenne de 35 roupies, roupie
valant 0.012~\$, fourchettes de 4~\% de la prime, frais de bourse et de régulation de 0.05~\% du
\omterm{def:m1:options-markets-europe-asia:turnover}{volume en primes}, taxe sur les transactions de 0.1~\% de la prime sur les ventes d'options. M2~: 850
millions de contrats, multiplicateur de 100, indice à 5\,600~\$, prime moyenne de 16~\$,
fourchettes de 1.5~\% de la prime, frais de 0.50~\$ par contrat pour un \omterm{def:m1:what-a-trading-firm-does:market-maker}{teneur de marché}.

\medskip\noindent\textbf{Partie I --- Trois mesures.}
\begin{enumerate}
\item Donner le rapport des contrats, M1 sur M2.
\item Donner le \omterm{def:m1:options-markets-europe-asia:turnover}{volume en notionnel} de chaque marché, en dollars.
\item Donner le \omterm{def:m1:options-markets-europe-asia:turnover}{volume en primes} de chaque marché, en dollars.
\item Donner le rapport des \omterm{def:m1:options-markets-europe-asia:turnover}{volumes en primes}, M2 sur M1.
\item Exprimer la prime de chaque marché en points de base du notionnel. Que dit
la différence sur ce qui est négocié~?
\end{enumerate}

\medskip\noindent\textbf{Partie II --- Un pool de revenus.}
\begin{enumerate}[resume]
\item Si les \omterm{def:m1:what-a-trading-firm-does:market-maker}{teneurs de marché} gagnent collectivement un quart de la fourchette
sur l'ensemble des primes, donner le pool brut de chaque marché.
\item Donner la taxe et les frais de M1 côté \omterm{def:m1:what-a-trading-firm-does:market-maker}{teneurs de marché}, en supposant que
les teneurs sont d'un côté de chaque transaction et vendent la moitié du temps.
\item Donner les frais de M2 pour les teneurs situés d'un côté de chaque transaction.
\item Donner les pools nets.
\item Quelle fraction de chaque pool un nouvel entrant pourrait-il réalistement
capter, et qu'est-ce qui la détermine~?
\end{enumerate}

\medskip\noindent\textbf{Partie III --- Risque réglementaire.}
\begin{enumerate}[resume]
\item Le régulateur de M1 triple la \omterm{def:m1:options-markets-europe-asia:turnover}{taille de lot} et supprime quatre \omterm{def:m1:options-markets-europe-asia:weekly}{échéances
hebdomadaires} sur cinq. Les contrats tombent à 24 milliards. Donner l'activité notionnelle conservée.
\item Donner le nouveau \omterm{def:m1:options-markets-europe-asia:turnover}{volume en primes} si la prime moyenne par unité reste
inchangée.
\item Recalculer le pool net de M1.
\item La société avait embauché douze personnes et loué une colocalisation pour
M1. Quelle leçon en tirer pour le business plan~?
\end{enumerate}

\medskip\noindent\textbf{Partie IV --- Jugement.}
\begin{enumerate}[resume]
\item Pourquoi le régulateur de M1 a-t-il agi~? Utiliser
l'\cref{dat:m1:options-markets-europe-asia:sebi}.
\item Qui se trouvait de l'autre côté des pertes des particuliers~?
\item Cela rend-il la tenue de marché sur M1 illégitime~? Où se situe la limite
que désigne la \cref{prop:m1:options-markets-europe-asia:leverage}~?
\item Qu'est-ce qu'une entreprise étrangère entrant sur M1 doit posséder, en plus d'une stratégie~?
\item Énoncer le \emph{résultat nommé}~: le rapport des \omterm{def:m1:options-markets-europe-asia:turnover}{volumes en primes}, M2 sur M1.
\item En une phrase~: que mesure un classement des contrats~?
\end{enumerate}
\end{problem}

\section{Questions d'entretien}

\begin{interviewq}[$\star$ \iqroles{trader, chercheur}]\label{iq:m1:options-markets-europe-asia:1}
Une bourse affirme être «~la plus grande bourse de produits dérivés au monde~». Que
demandez-vous~?
\end{interviewq}

\begin{interviewq}[$\star$ \iqroles{trader, chercheur, banque}]\label{iq:m1:options-markets-europe-asia:2}
En quoi les marchés d'options cotées en Europe diffèrent-ils de ceux des
États-Unis~?
\end{interviewq}

\begin{interviewq}[$\star\star$ \iqroles{chercheur, trader}]\label{iq:m1:options-markets-europe-asia:3}
Un régulateur triple la taille minimale de contrat des options sur indice. Prédire
les effets sur le volume, sur les participants et sur les revenus des \omterm{def:m1:what-a-trading-firm-does:market-maker}{teneurs de marché}.
\end{interviewq}

\begin{interviewq}[$\star\star$ \iqroles{trader, chercheur}]\label{iq:m1:options-markets-europe-asia:4}
Pourquoi les produits dérivés sur indice \omterm{def:m1:basis-roll-and-delivery:delivery}{réglés en espèces} sont-ils vulnérables
autour de leur règlement-livraison, et comment les bourses s'en protègent-elles~?
\end{interviewq}

\begin{interviewq}[$\star\star$ \iqroles{trader, banque}]\label{iq:m1:options-markets-europe-asia:5}
Votre desk au comptant et votre desk d'options sont tous deux actifs sur le même
indice le jour de l'échéance. Quels contrôles souhaitez-vous mettre en place~?
\end{interviewq}

\begin{interviewq}[$\star\star\star$ \iqroles{chercheur, trader}]\label{iq:m1:options-markets-europe-asia:6}
On vous demande d'évaluer l'entrée sur un marché d'options de détail en forte
croissance à l'étranger. Structurez l'analyse.
\end{interviewq}
